\section{Literature review}
\label{sec:literature review}


The problem of chess puzzle generation has been studied quite extensively despite its limited practical applicability \cite{feng2025generatingcreativechesspuzzles, iqbal2021acomputationalmethodofoptimizingchesscompositions, iqbal2017computercomposesfabledproblem}. Chess platforms such as Chess.com and Lichess primarily generate puzzles by filtering large databases of human games to find positions with unique solutions. In the scientific community, more effort has been directed towards procedurally generating chess positions with desired properties.

One traditional approach to chess puzzle generation employs evolutionary search algorithms \cite{feng2025generatingcreativechesspuzzles}. This process typically begins with a random board state and iteratively applies minor random modifications to improve the position. A fitness function evaluates each configuration, advancing the most promising candidates until a satisfactory puzzle emerges. \citeauthor{feng2025generatingcreativechesspuzzles} \cite{feng2025generatingcreativechesspuzzles} briefly discuss their initial trials with this technique, leaving a comprehensive exploration for future work. Although their primary methodology relies on deep learning, they suggest that a hybrid approach incorporating evolutionary search could be a good place for future research. Alternatively, one can employ an iterative refinement approach. Unlike evolutionary algorithms that maintain a large population of diverse positions, this method focuses on a single, pre-existing puzzle, incrementally applying small, deterministic modifications to optimize its aesthetic qualities \cite{iqbal2021acomputationalmethodofoptimizingchesscompositions, iqbal2017computercomposesfabledproblem}. Notably, because the starting position is generally assumed to be a valid puzzle already, this functions more as a refinement technique than true generation.

In the field of chess puzzle generation, \citeauthor{feng2025generatingcreativechesspuzzles} \cite{feng2025generatingcreativechesspuzzles} demonstrate the viability of using deep learning models, in particular language models for chess puzzle generation. They train four language models with different architectures using supervised training and then fine-tune the autoregressive model further using reinforcement learning. Their goal was to generate positions that humans would consider creative. They found that reinforcement learning substantially improved both the quality and quantity of the generated positions, ultimately producing positions comparable to award-winning compositions evaluated by master-level players and puzzle composers.

While autoregressive models have proven capable in this domain, the potential of diffusion large language models (dLLMs) remain comparably unexplored. Masked diffusion models, originally developed by \citeauthor{austin2021structureddenoisingdiffusionmodels} \cite{austin2021structureddenoisingdiffusionmodels} and later scaled up by others (see e.g., \cite{nie2025largelanguagediffusionmodels, ye2025dream7bdiffusionlarge}), offer tradeoffs over autoregressive architectures. In particular, they offer benefits in terms of efficiency during inference, as one can vary the amount of inference time compute depending on the task by using a different number of denoising steps. In addition, they are able to use bidirectional attention, which may offer benefits compared to autoregressive models. However, they take much longer to train \cite{nie2025scalingmaskeddiffusionmodels}.

A critical component of modern language model training is reinforcement learning, such as Reinforcement Learning from Human Feedback (RLHF) and Reinforcement Learning with Verifiable Rewards (RLVR) \cite{shao2024deepseekmathpushinglimitsmathematical, ouyang2022traininglanguagemodelstofollowinstructionswithhumanfeedback, lambert2025tulu}. However, its application to discrete diffusion models is not as thoroughly studied. While \citeauthor{black2024trainingdiffusionmodelsreinforcement} \cite{black2024trainingdiffusionmodelsreinforcement} provided a framework for applying reinforcement learning to continuous diffusion models, extending these methods to language modeling presents computational challenges due to the difficulty of evaluating exact token-wise transition probabilities within one forward pass, which is possible for autoregressive LLMs. Therefore, new methods have been proposed to mimic the work of \citeauthor{zheng2025groupsequencepolicyoptimization} \cite{zheng2025groupsequencepolicyoptimization} and group sequence policy optimization (GSPO).

Recent works have proposed methods to approximate the likelihood of individual token sequences in dLLMs \cite{ou2025principledrldiffusionllms, zhao2025d1scalingreasoningdiffusion, tang2026wd1weightedpolicyoptimization, rojas2025improvingreasoningdiffusionlanguage}, paving the way for RL techniques similar to those used in standard LLMs. For instance, ELBO-based sequence-level policy optimization (ESPO) \cite{ou2025principledrldiffusionllms} has been successfully applied to dLLM post-training. Compared to Denoising Diffusion Policy Optimization (DDPO) \cite{black2024trainingdiffusionmodelsreinforcement}, ESPO works with a sequence-level Markov Decision Process (MDP), which allows one to use the evidence lower bound as a surrogate for the log probability of the sample. Instead, DDPO works with a step-wise MDP and exact transition probabilities offering unbiased gradient updates. Although DDPO is computationally less efficient than ESPO, meaning that the model update for each sampled trajectory takes quite a bit longer with DDPO, than with ESPO, DDPO allows the model to make updates for each sampled token individually (in reality DDPO works with all tokens sampled on a given denoising step, but this is still better than the whole sequence). This step-wise update could offer benefits in tasks, where individual tokens matter a lot. In summary, while RL has proven effective for autoregressive chess generation, its application to dLLMs remains an open challenge. To address this, we explore adapting DDPO for masked diffusion models, leveraging step-wise token updates to generate and optimize high-quality chess puzzles.
